\documentclass[11pt]{article}
\usepackage[T1]{fontenc}
\usepackage[a4paper,margin=27mm]{geometry}
\usepackage{lmodern,amsmath,amssymb,amsthm,microtype,hyperref}
\hypersetup{hidelinks}
\newtheorem{theorem}{Theorem}[section]
\newtheorem{lemma}[theorem]{Lemma}
\newtheorem{proposition}[theorem]{Proposition}
\newcommand{\N}{\mathrm N}
\newcommand{\Q}{\mathbb Q}
\newcommand{\workbenchurl}{https://github.com/KokunoYumeto/zeta-and-l-functions-workbench}
\newcommand{\recordurl}{https://github.com/KokunoYumeto/mathematics-workbenches}
\title{A logarithmic convolution criterion for cubic RH, after Hecke}
\author{}
\date{24 November 2026}
\begin{document}
\maketitle
\begin{abstract}
Hecke's functional equation controls smooth sums of a fixed character, and Gevrey's derivative classes make the dependence on the smoothing explicit. For the cubic character defining $\mathbb Q(\sqrt[3]{269})$, we reduce the smooth prime sum to an exact convolution whose direct character variable is at most a constant times $(\log X)^2(\log\log X)^\kappa$, for any fixed $\kappa>2$. The discarded part is smaller than every negative power of $X$. A further arithmetic decomposition removes rational contents at least $Y^{1/3}$ from each inverse norm window $Y$, at square-root cost; classical zeta RH improves this threshold to $Y^{2/7}$. The test functions are fixed, nonnegative, and can be supported in any prescribed subinterval of the dyadic interval. A square-root bound for the retained inverse range is equivalent to the cubic Riemann hypothesis and is not proved here. We give full proofs of the cutoff and content estimates, explain the separate-estimation loss, and show that smooth numerator truncation and the primitive inverse series preserve the relevant zero poles. The mathematics and text were produced by GPT-6 Astra in ChatGPT Pro mode; no human has checked this text.
\end{abstract}

\section{The precise reduction and its scope}
We state an unconditional reduction for one fixed cubic character.
Let $K=\Q(\omega)$, where $\omega^2+\omega+1=0$, and let $\psi$
be either nontrivial character of
$\operatorname{Gal}(K(\sqrt[3]{269})/K)$, viewed as a primitive
finite-order Hecke character. Its Hecke function is entire and
\[
 L_K(s,\psi)=
 \frac{\zeta_{\Q(\sqrt[3]{269})}(s)}{\zeta(s)}.
\]
Hecke's theorem, in the adelic proof of Tate, gives the functional
equation and polynomial growth in every closed vertical strip.
Poonen's full exposition proves these analytic statements in
Theorems~5.16 and 5.22(b,c), pp.~34--39.
Let $S$ consist of the prime ideals above 3 and 269, and put
\[
 P(s)=L_K^S(s,\psi),\qquad
 a(\mathfrak c)=\psi(\mathfrak c),\qquad
 m(\mathfrak d)=\mu_K(\mathfrak d)\psi(\mathfrak d).
\]
All ideals in what follows are nonzero integral ideals coprime to $S$.
Here $\mu_K$ is the ideal M\"obius function, and $\N$ is the absolute
ideal norm. The finite Euler deletions leave $P$ entire and preserve
the stated strip growth.
Define $\Lambda_K(\mathfrak p^k)=\log\N\mathfrak p$ for prime ideals
$\mathfrak p$ and integers $k\ge1$, and define it to be zero on all
other ideals.

Fix $W\in C_c^\infty((1,2))$ and a smooth real function $\chi$ on
$[0,\infty)$ with $0\le\chi\le1$, $\chi(v)=1$ for $v\le1$,
and $\chi(v)=0$ for $v\ge2$. Put
\[
 T_W(X)=\sum_{\mathfrak n}
       \Lambda_K(\mathfrak n)\psi(\mathfrak n)W(\N\mathfrak n/X).
\]
For $0<\delta<1$, write $D=X^{1-\delta}$ and define the exact sum
\begin{align}
 R_\delta(X)
 &=\sum_{\N\mathfrak c<2X^\delta}
       a(\mathfrak c)\log\N\mathfrak c
       \sum_{\mathfrak d}m(\mathfrak d)
       W(\N\mathfrak c\N\mathfrak d/X)
       \bigl(1-\chi(\N\mathfrak d/D)\bigr).
                                                        \label{eq:remainder}
\end{align}

\begin{theorem}\label{thm:switch}
For the fixed character, fixed weights, and every fixed
$0<\delta<1$ and $M>0$,
\[
 T_W(X)=R_\delta(X)+O_{W,\chi,\delta,M}(X^{-M}),
       \qquad X\ge2.                                      \tag{1.1}
\]
This assertion is unconditional. If, in addition, $P$ has no zero
in $\Re s>\beta$, where $1/2\le\beta<1$, then separate estimation
of the inner M\"obius sums gives, for every $\epsilon>0$,
\[
 R_\delta(X)\ll_{W,\chi,\delta,\beta,\epsilon}
 X^{\beta+\delta(1-\beta)+\epsilon}.                        \tag{1.2}
\]
The complete sum has the stronger bound
$T_W(X)=O_{W,\beta}(X^\beta)$ under the same zero-location
hypothesis, so (1.1) also gives $R_\delta(X)=O_{W,\chi,\delta,\beta}(X^\beta)$.
No square-root bound for $R_\delta$ is asserted.
\end{theorem}

The constants in (1.1) depend on the fixed positive $\delta$.
The statement does not allow $\delta$ to tend to zero with $X$.
By itself this estimate does not make the direct variable
$\N\mathfrak c$ logarithmic in $X$. The stronger derivative control
in the next theorem permits such a conclusion with explicit uniformity.
For the prime terms, deleting prime powers costs $O_W(\sqrt X)$,
so a square-root bound for the displayed remainder would give
the corresponding square-root bound for the cubic prime sum.
The theorem supplies no such estimate.

\subsection{The logarithmic cutoff criterion}
Write $\mathcal D=u\,d/du$. For any $L>1$, set
\begin{align*}
 R_W(X;L)={}&\sum_{\N\mathfrak c<2L}
     \psi(\mathfrak c)\log\N\mathfrak c\\
 &\quad\times\sum_{\mathfrak d}\mu_K(\mathfrak d)\psi(\mathfrak d)
 W(\N\mathfrak c\N\mathfrak d/X)
 \bigl(1-\chi(L\N\mathfrak d/X)\bigr).
\end{align*}
This is a finite sum, and the nonzero inverse terms satisfy
$\N\mathfrak d>X/L$.

\begin{theorem}\label{thm:near-analytic-cutoff}
Fix $r>1$ and $\kappa>2r$. Let the fixed function
$W\in C_c^\infty((1,2))$ satisfy, for fixed constants $C,H>0$,
\[
 \|\mathcal D^kW\|_\infty\le C H^k k![\log(e+k)]^{rk}
 \qquad(k\ge0).
\]
Put
\[
 \ell_1(X)=\log(3X),\quad
 \ell_2(X)=\log(e+\log(3X)),\quad
 L_X=\ell_1(X)^2\ell_2(X)^\kappa.
\]
For every $M>0$,
\[
 T_W(X)=R_W(X;L_X)+O_{W,\chi,r,\kappa,M}(X^{-M})
 \qquad(X\ge2).
\]
Nonzero nonnegative functions with the stated derivative bound exist
in every nonempty open subinterval of $(1,2)$. Thus every fixed
$\kappa>2$ is allowed by first choosing $1<r<\kappa/2$.
For these fixed parameters, RH for the primitive cubic Hecke
function is equivalent to
\[
 R_W(X;L_X)=O_{W,\epsilon}(X^{1/2+\epsilon})
\]
for every $\epsilon>0$ and every such nonnegative test function.
This bound is not asserted.

If every nontrivial zero has real part at most
$\beta$, $1/2\le\beta<1$, the complete smooth explicit formula
gives $R_W(X;L_X)=O_W(X^\beta)$. Bounding its inner M\"obius
sums separately gives
\[
 R_W(X;L_X)\ll_{W,\epsilon} X^{\beta+\epsilon}
                    L_X^{1-\beta}\log(2L_X).
\]
All implicit constants may also depend on the fixed character,
cutoff, and displayed fixed parameters.
\end{theorem}
The theorem reduces the support of the remaining convolution;
it does not improve the known zero-free exponent. Its full proof
is given in Sections~\ref{sec:gevrey} and~\ref{sec:near-analytic}.
The endpoint $\kappa=2$ and arbitrary smooth weights are not included.
The classical derivative classes are due to Gevrey
(1918, Part I, Section 1, p.~132); the estimates and compact-weight
construction needed here are proved in full.

\subsection{Removing large rational content}
For an Eisenstein ideal $\mathfrak a=(a+b\omega)$ put
$\operatorname{cont}(\mathfrak a)=\gcd(a,b)>0$.
The Eisenstein ring is a principal ideal domain, and its units
preserve this gcd. Thus the definition is intrinsic.
For $Y_{\mathfrak c}=X/\N\mathfrak c$, let
$R^{\rm low}_{W,\gamma}(X;L)$ be the sum $R_W(X;L)$ with the
additional inner multiplier
\[
 \chi\left(\operatorname{cont}(\mathfrak d)/
                          Y_{\mathfrak c}^{\gamma}\right).
\]

\begin{theorem}\label{thm:content-reduction}
Use the weights, parameters, and logarithmic length $L_X$ of
Theorem~\ref{thm:near-analytic-cutoff}.
For every $\epsilon>0$,
\[
 T_W(X)=R^{\rm low}_{W,1/3}(X;L_X)
                         +O_{W,\chi,\epsilon}(X^{1/2+\epsilon}).
\]
If classical RH for the Riemann zeta function is granted, then
\[
 T_W(X)=R^{\rm low}_{W,2/7}(X;L_X)
                         +O_{W,\chi,\epsilon}(X^{1/2+\epsilon}).
\]
Consequently cubic RH is equivalent to a square-root bound
for $R^{\rm low}_{W,1/3}$ for every positive error exponent and
every nonnegative localized weight in the stated class.
Under zeta RH the same equivalence holds with $1/3$ replaced
by $2/7$.

In the retained sum, writing
$\mathfrak d=(u\alpha_0)$ with $u=\operatorname{cont}(\mathfrak d)$
and $\alpha_0$ primitive, one has
\[
 \N\mathfrak c<2L_X,\quad
 Y_{\mathfrak c}<\N\mathfrak d<2Y_{\mathfrak c},\quad
 u<2Y_{\mathfrak c}^{\gamma},\quad
 \N\alpha_0>\tfrac14Y_{\mathfrak c}^{1-2\gamma}.
\]
The theorem bounds the discarded content range only.
The square-root estimate for either retained sum is not proved.
\end{theorem}

Section~\ref{sec:content} proves the arithmetic estimates
\[
 \sum_{\mathfrak a}\mu_K(\mathfrak a)\psi(\mathfrak a)
 V(\N\mathfrak a/Y)
 (1-\chi(\operatorname{cont}(\mathfrak a)/R))
 \ll_\epsilon Y^{1+\epsilon}R^{-3/2},
\]
and, under zeta RH, the same bound with $R^{-7/4}$.
These improve the absolute content-tail bound $O(Y/R)$.
They delete a sparse inverse range at square-root cost,
while leaving the primitive component $u=1$.


\subsection{Comparison with the squarefree factor}

There is a useful comparison with the squarefree factor.
The coefficients of $B(s)=P(s)/P(2s)$ are
$b(\mathfrak c)=\mu_K(\mathfrak c)^2\psi(\mathfrak c)$, and its inverse
has coefficients $(-1)^{\Omega(\mathfrak d)}\psi(\mathfrak d)$.
The denominator $P(2s)$ permits a smooth coefficient estimate only
at exponent $\beta/2+\epsilon$ from the asserted zero-free region.
Absolute summation of that inverse up to $D$ gives
$X^{\beta/2+\epsilon}D^{1-\beta/2}\log(2X)$, reaching the
square-root scale at $D=X^{(1-\beta)/(2-\beta)}$.
In contrast, the numerator $P$ itself is entire, which allows
the unconditional deletion in (1.1). The price is the
M\"obius coefficients in the long sum \eqref{eq:remainder}.
This is a reduction of the remaining estimate, not a proof of RH.

The fixed-conductor rapid decrease used below is classical.
OpenAI (2026a, Lemma~4.2, p.~14) gives a conductor-dependent
full-factor reflection with dual length proportional to the conductor
divided by the original length. The same consequence when that dual
length is smaller than one is explicitly recorded by
bramble\_trace (2026, Section~2.2).
We write its elementary Mellin proof to specify exactly what
uniformity survives the summation over the inverse variable.
For numerical comparison of the remaining exponent, OpenAI's
September 30 Theorem~1.1 gives $\beta=7/8$.
Liu's main Theorems~1.2 and 8.29 give the sharper boundary
\[
 B_*=0.873345856106718576801184796065\ldots.
\]
This is the degree-nine algebraic constant of his equations (42)
and (44), specified by the root in
$(0.8733458561067185768,0.8733458561067185769)$.
His rational one-stage value $3493459/4000000$ is a separate
intermediate result.
These zero-free estimates are external inputs used only for the numerical
comparison of the remaining exponent. The unconditional convolution
reductions use neither.

\section*{Proof overview}
The direct factor has no pole: its logarithmically weighted
coefficient sum is a Mellin integral of $-P'$, an entire function.
Moving that integral arbitrarily far to the left gives a negative
power of the direct length. When the inverse variable is at most
$X^{1-\delta}$, the direct length is at least $X^\delta$.
A sufficiently large fixed number of contour shifts therefore
absorbs the entire number of inverse coefficients.

Switching the remaining finite convolution forces the direct variable
below $2X^\delta$. A smooth cutoff makes the inner weights range
over a bounded family, so the inverse estimate can be applied uniformly.
For the stronger cutoff, we move the contour an integer distance $A$ and retain its dependence on $A$. The degree-two gamma quotient costs essentially $(2A)!$. Controlled derivatives permit optimizing $A$ with the scale. Summable interval widths produce nearly analytic compact weights and the logarithmic direct range. Localized nonnegative weights then detect any proposed zero pole through their Mellin transforms.

The content decomposition writes a nonzero inverse coefficient as
\[
 \mu(u)^2\chi_{-3}(u)\mu(\N\alpha_0)\psi((\alpha_0)),
\]
with the required coprimality condition. Summing in the rational
content $u$ first gives the additional support reduction.
Its elementary bound uses only the squarefree expansion and
periodicity of $\chi_{-3}$; the stronger smooth bound uses zeta RH.

The zero-free region controls $1/P$ only to the right of $\beta$;
summing the direct variable in absolute value loses
$\delta(1-\beta)$. Finally, a contour around a zero shows that
smoothly truncating the direct factor preserves its nonzero
logarithmic-derivative residue. Entireness of the direct factor
therefore does not cancel the zero poles.

\section{Rapid decrease and the switched remainder}
We first prove the fixed-character cancellation with the strength
needed for a polynomial number of inverse coefficients.

\begin{lemma}\label{lem:rapid}
For every fixed $A>0$,
\[
 A_W(Y):=\sum_{\mathfrak c}a(\mathfrak c)\log\N\mathfrak c\,
 W(\N\mathfrak c/Y)=O_{W,A}(Y^{-A}),\qquad Y\ge1.             \tag{2.1}
\]
\end{lemma}
\begin{proof}
For $c>1$, absolute convergence gives Mellin inversion
\[
 A_W(Y)=\frac1{2\pi i}\int_{\Re s=c}
              -P'(s)\widehat W(s)Y^s\,ds,\qquad
 \widehat W(s)=\int_0^\infty W(u)u^{s-1}\,du.
\]
The transform is entire and rapidly decreasing on each closed
vertical strip. The fixed Hecke functional equation, Stirling's
formula, and absolute convergence on the right give polynomial
growth of $P$ on any such strip; Cauchy's formula gives the
same statement for $P'$. Thus the contour moves to $\Re s=-A$
without residues. The integral there is absolutely convergent
and is $O_{W,A}(Y^{-A})$. All constants may depend on the
fixed character and its finite Euler deletions.
\end{proof}

\begin{proof}[Proof of Theorem~\ref{thm:switch}, unconditional assertion]
The absolutely convergent Euler products on $\Re s>1$ give
\[
 \frac1{P(s)}=\sum_{\mathfrak d}
                  \frac{m(\mathfrak d)}{(\N\mathfrak d)^s},
 \qquad
 -\frac{P'(s)}{P(s)}
  =\sum_{\mathfrak n}
     \frac{\Lambda_K(\mathfrak n)\psi(\mathfrak n)}
          {(\N\mathfrak n)^s}.
\]
Coefficient comparison gives the exact convolution
\[
 T_W(X)=\sum_{\mathfrak d}m(\mathfrak d)
            A_W(X/\N\mathfrak d).
                                                               \tag{2.2}
\]
It is a finite sum because $W$ has compact support.
The part multiplied by $\chi(\N\mathfrak d/D)$ has
$\N\mathfrak d\le2D$. Lemma~\ref{lem:rapid} and
$|m(\mathfrak d)|\le1$ show that, for sufficiently large $X$,
its absolute value is at most
\[
 C_{W,A}X^{-A}
   \sum_{\N\mathfrak d\le2D}(\N\mathfrak d)^A
 \ll_{W,A}X^{-A}D^{A+1}.
                                                               \tag{2.3}
\]
The ideal count in the fixed field is $O(t)$, which also follows
directly from its Eisenstein norm lattice. This justifies the last
estimate. With $D=X^{1-\delta}$, the exponent in (2.3) is
$1-\delta-A\delta$. Choose
$A>(M+1-\delta)/\delta$. Bounded $X$ is absorbed in the constant.
The other part of (2.2) has $\N\mathfrak d>D$ and
$\N\mathfrak c\N\mathfrak d<2X$, so $\N\mathfrak c<2X^\delta$.
Interchanging the two finite sums gives exactly
\eqref{eq:remainder}. This proves (1.1).
\end{proof}

\begin{proof}[Proof of the remaining bounds]
In a closed inner strip of $\Re s>\beta$, $1/P(s)$ has
polynomial growth. To see the dependence needed here, choose an
analytic logarithm in that half-plane, anchored by the bounded Euler
logarithm at $2+it$. A disk centered at that point and lying in a
slightly larger inner strip reaches any prescribed point of the
smaller strip. Polynomial growth gives
$\log|P|=O(\log(2+|t|))$ from above on the disk.
Borel--Carath\'eodory gives
$|\log P|=O(\log(2+|t|))$ on the smaller disk; exponentiating
gives the claimed reciprocal bound. For $\Re s\ge2$ the Euler
product gives it directly.
Consequently Mellin inversion of $1/P$ gives
\[
 \sum_{\mathfrak d}m(\mathfrak d)V(\N\mathfrak d/Y)
       \ll_{\beta,\epsilon,V}Y^{\beta+\epsilon},\qquad Y\ge1,
                                                               \tag{2.4}
\]
for every smooth compactly supported $V$.
The bound is uniform for weights with common compact support
and uniformly bounded derivatives of sufficiently many orders.

In \eqref{eq:remainder} set $Y=X/\N\mathfrak c$ and
$\lambda=Y/D$. The inner weight in (2.4) is
\[
 V_\lambda(u)=W(u)\bigl(1-\chi(\lambda u)\bigr).
\]
The outer range has $\lambda>1/2$. If $\lambda\ge2$ then
$V_\lambda=W$. For $1/2<\lambda<2$ these weights have uniformly
bounded derivatives and the common support of $W$.
Thus (2.4) applies uniformly throughout the outer sum.
Choose $0<\epsilon'<1-\beta$. Counting ideals gives
\begin{align*}
 |R_\delta(X)|
 &\ll X^{\beta+\epsilon'}
    \sum_{\N\mathfrak c<2X^\delta}
       \frac{\log\N\mathfrak c}{(\N\mathfrak c)^{\beta+\epsilon'}}\\
 &\ll X^{\beta+\delta(1-\beta)+(1-\delta)\epsilon'}\log(2X).
\end{align*}
Taking $\epsilon'$ sufficiently small and absorbing the logarithm
proves (1.2).

For completeness, the stronger bound for the complete sum follows
from the usual smooth explicit formula for $-P'/P$.
There is no pole at one. The fixed-conductor local zero count is
$O(\log(2+|t|))$; see Greni\'e and Molteni (2019, Appendix~A,
Theorem~A.1). The Mellin transform decreases faster than any
power of $|t|$, and the asserted zero location bounds each
nontrivial-zero term by $X^\beta$ times that transform.
The trivial zeros and the finitely many deleted local factors
cost $O_W(\log(2X))$, which is absorbed by $X^\beta$.
Thus $T_W(X)=O_{W,\beta}(X^\beta)$.
Equation (1.1) transfers this bound to $R_\delta$.
\end{proof}

For $\beta=7/8$, separate inner estimation gives
$7/8+\delta/8+\epsilon$, and the whole-sum estimate remains $7/8$.
For Liu's main theorem the separate exponent is
$B_*+\delta(1-B_*)+\epsilon$, and the complete-sum exponent is
$B_*$. Neither becomes $1/2$ by changing $\delta$.

\section{Logarithmic direct length with controlled derivatives}
\label{sec:gevrey}
The functional equation gives a stronger cutoff when the derivatives of
our fixed test function have specified factorial growth. This section
keeps track of that growth before allowing the cutoff to depend on the
scale. The class of functions is the classical Gevrey class; we prove
below the existence of every nonnegative test function that the argument
needs.

Write $\mathcal D=u\,d/du$. Fix $\sigma>1$, and suppose that
$W\in C_c^\infty((1,2))$ satisfies
\begin{equation}\label{eq:gevrey-condition}
 \|\mathcal D^k W\|_\infty\le C_W H_W^k(k!)^\sigma
 \qquad(k=0,1,2,\ldots).
\end{equation}
The constants $C_W,H_W$ are fixed independently of the scale.

\begin{theorem}\label{thm:gevrey-decay}
For the fixed cubic character, there exist constants $C,c>0$ such that
\begin{equation}\label{eq:gevrey-decay}
 A_W(Y)=\sum_{\mathfrak c}\psi(\mathfrak c)\log\N\mathfrak c\,
 W(\N\mathfrak c/Y)
 \quad\hbox{satisfies}\quad
 |A_W(Y)|\le C\exp\bigl(-cY^{1/(2\sigma)}\bigr)
 \qquad(Y\ge1).
\end{equation}
These constants may depend on the fixed character, the omitted Euler
factors, and the constants in \eqref{eq:gevrey-condition}.
\end{theorem}

\begin{proof}
Put $q=217083$ and $a_0=2\pi/\sqrt q$. The primitive Hecke function has
completion $2a_0^{-s}\Gamma(s)L(s)$ and functional equation relating
$s$ to $1-s$. Its dual has an absolutely convergent Dirichlet series
on every half-plane $\Re s\ge1+\epsilon$. We first obtain a bound whose
dependence on the distance moved to the left is explicit.

For an integer $A\ge1$, write $s=-A+z$. The gamma recurrence gives
\[
 \frac{\Gamma(1-s)}{\Gamma(s)}
 =(-1)^A\prod_{j=1}^A(j-z)^2
          \frac{\Gamma(1-z)}{\Gamma(z)}.
\]
On $|\Re z|\le1/4$, Stirling's estimate and boundedness on compact
sets give
\[
 \left|\frac{\Gamma(1-z)}{\Gamma(z)}\right|
 \le C(1+|\Im z|)^{3/2}.
\]
There is no pole of this quotient in that strip; the apparent point
$z=0$ is a zero. If $|z-it|\le1/4$, then
$|j-z|\le A+|t|+1/2\le(3/2)(A+|t|)$.
The functional equation, absolute convergence of the dual series
on $\Re(1-s)\ge7/4$, and the finite number of deleted Euler factors
therefore give
\[
 \sup_{|s-(-A+it)|\le1/4}|P(s)|
 \le C_0 C_1^A(A+|t|)^{2A+2}.
\]
Here and below $C_0,C_1$ do not depend on $A$ or $t$. The conductor
factor contributes a fixed constant to the power $A$; each deleted
Euler factor does the same. Cauchy's derivative formula on the circle
of radius $1/4$ gives
\begin{equation}\label{eq:gevrey-left-growth}
 |P'(-A+it)|\le C_2 C_3^A(A+|t|)^{2A+2}.
\end{equation}

For $s=-A+it$, integrate the Mellin transform by parts $k$ times
using $\mathcal D$. Boundary terms vanish because $W$ is supported
in a compact subinterval of $(1,2)$, and $u^{-A}\le1$ there. Thus
\begin{equation}\label{eq:gevrey-mellin}
 |\widehat W(-A+it)|
 \le C_W(\log2)H_W^k(k!)^\sigma(A^2+t^2)^{-k/2}.
\end{equation}
The ordinary fixed-contour argument first gives
\[
 A_W(Y)=\frac1{2\pi i}\int_{\Re s=-A}
                   -P'(s)\widehat W(s)Y^s\,ds
\]
for each integer $A$. No pole is crossed, since $P'$ is entire.
We may choose a different integer $A$ for each $Y$, because the
identity holds for all such integers and the estimates now have
explicit uniform dependence on $A$.

Set $k=2A+4$. Since $A+|t|\le\sqrt2(A^2+t^2)^{1/2}$, equations
\eqref{eq:gevrey-left-growth} and \eqref{eq:gevrey-mellin} imply
\begin{align}
 |A_W(Y)|
 &\le C_4 C_5^A((2A+4)!)^\sigma Y^{-A}
          \int_{\mathbb R}\frac{dt}{A^2+t^2}\notag\\
 &\le C_6 C_5^A((2A+4)!)^\sigma Y^{-A}.
                                                        \label{eq:gevrey-factorial}
\end{align}
Choose $A=\lfloor c_0Y^{1/(2\sigma)}\rfloor$, with $c_0>0$
sufficiently small. The elementary bound $k!\le k^k$ shows that
the logarithm of the right side of \eqref{eq:gevrey-factorial}
is at most $-c_1Y^{1/(2\sigma)}$ once $Y$ is large: the term
$2\sigma A\log A$ cancels $A\log Y$, while the small fixed
$c_0$ leaves a negative constant multiple of $A$; the extra
$4\sigma\log(2A+4)$ is absorbed. Increasing $C$ handles bounded
$Y$ and proves \eqref{eq:gevrey-decay}.
\end{proof}

For a fixed $\kappa>2\sigma$, put
\[
 L_X=(\log(3X))^\kappa,\qquad D_X=X/L_X.
\]
Define
\begin{align}
 R_W(X;L_X)
 ={}&\sum_{\N\mathfrak c<2L_X}
       \psi(\mathfrak c)\log\N\mathfrak c\notag\\
 &\quad\times\sum_{\mathfrak d}\mu_K(\mathfrak d)\psi(\mathfrak d)
       W(\N\mathfrak c\N\mathfrak d/X)
       \bigl(1-\chi(\N\mathfrak d/D_X)\bigr).
                                                       \label{eq:log-remainder}
\end{align}
The inverse variable is still long: it exceeds $X/L_X$ wherever
the cutoff is nonzero. The direct character variable is logarithmic.

\begin{theorem}\label{thm:log-cutoff}
For every $M>0$,
\begin{equation}\label{eq:log-cutoff}
 T_W(X)=R_W(X;L_X)+O_{W,\chi,\sigma,\kappa,M}(X^{-M})
 \qquad(X\ge2).
\end{equation}
This assertion is unconditional and permits every fixed $\kappa>2$
by first choosing $1<\sigma<\kappa/2$ and a test function satisfying
\eqref{eq:gevrey-condition}. It does not assert the endpoint
$\kappa=2$.
\end{theorem}

\begin{proof}
Use the exact finite convolution
$T_W(X)=\sum_{\mathfrak d}\mu_K(\mathfrak d)\psi(\mathfrak d)
A_W(X/\N\mathfrak d)$.
The part with cutoff $\chi(\N\mathfrak d/D_X)$ has
$\N\mathfrak d\le2D_X$, so $X/\N\mathfrak d\ge L_X/2$.
There are $O(D_X)$ such ideals when the range is nonempty;
this is the elementary Eisenstein norm-lattice count. Theorem
\ref{thm:gevrey-decay} bounds the discarded part by
\[
 C X\exp\bigl(-c(L_X/2)^{1/(2\sigma)}\bigr).
\]
Since $\kappa/(2\sigma)>1$, this is $O_M(X^{-M})$ for every
fixed $M$. In the remaining sum $\N\mathfrak d>D_X$ and
$\N\mathfrak c\N\mathfrak d<2X$, which force
$\N\mathfrak c<2L_X$. Interchanging the finite sums gives
\eqref{eq:log-remainder}. Bounded $X$ is absorbed in the constant.
\end{proof}

\begin{lemma}\label{lem:gevrey-bumps}
For every $\sigma>1$ and every nonempty open interval
$J\subset(1,2)$, there is a nonzero nonnegative function
$W\in C_c^\infty(J)$ satisfying \eqref{eq:gevrey-condition}.
\end{lemma}

\begin{proof}
Put $r=1/(\sigma-1)$ and define $h(x)=\exp(-x^{-r})$ for $x>0$
and $h(x)=0$ for $x\le0$. Choose a fixed $\eta>0$, depending
only on $r$, so small that on $|z-x|\le\eta x$, $x>0$, the
principal branch satisfies $\Re(z^{-r})\ge c_rx^{-r}$.
Cauchy's formula gives
\[
 |h^{(k)}(x)|\le k!(\eta x)^{-k}\exp(-c_rx^{-r}).
\]
For $k\ge1$, maximizing $v^{k/r}e^{-c_rv}$ over $v>0$ and
using $k!\ge(k/e)^k$ bounds this by
$C_r^k(k!)^{1+1/r}=C_r^k(k!)^\sigma$. For each fixed $k$,
the same estimate before maximization tends to zero as $x\downarrow0$,
so all derivatives extend continuously by zero and $h$ is smooth.
The case $k=0$ is bounded directly.

Choose $a<b$ with $[a,b]\subset\log J$. The function
$F(v)=h(v-a)h(b-v)$ is nonnegative, supported on $[a,b]$,
and positive inside. Leibniz's formula and
\[
 \binom{k}{j}(j!)^\sigma((k-j)!)^\sigma
 =(k!)^\sigma \binom{k}{j}^{1-\sigma}\le(k!)^\sigma
\]
give $\|F^{(k)}\|_\infty\le C H^k(k!)^\sigma$, after
absorbing $k+1$ into $H^k$. Take $W(u)=F(\log u)$.
Then $\mathcal D^kW(u)=F^{(k)}(\log u)$, proving the lemma.
\end{proof}

\begin{proposition}\label{prop:gevrey-rh}
Fix $\kappa>2$ and $1<\sigma<\kappa/2$. The Riemann hypothesis
for the primitive cubic Hecke function is equivalent to
\begin{equation}\label{eq:gevrey-rh}
 R_W(X;L_X)=O_{W,\epsilon}(X^{1/2+\epsilon})
 \quad\hbox{for every }\epsilon>0
\end{equation}
for every nonnegative $W$ satisfying \eqref{eq:gevrey-condition}.
Under a zero-free half-plane $\Re s>\beta$, with
$1/2\le\beta<1$, the complete smooth explicit formula gives
$R_W(X;L_X)=O_W(X^\beta)$. Separately bounding the inner
M\"obius sums gives only
\begin{equation}\label{eq:gevrey-separate}
 R_W(X;L_X)
 \ll_{W,\epsilon}X^{\beta+\epsilon}
        (\log(3X))^{\kappa(1-\beta)}\log\log(9X).
\end{equation}
\end{proposition}

\begin{proof}
The whole-sum explicit formula proved above bounds $T_W(X)$ by
$O_W(X^\beta)$ if all its nontrivial zeros have real part at most
$\beta$. Theorem~\ref{thm:log-cutoff} transfers that bound to the
remainder, proving the forward RH implication as well.

Conversely, suppose \eqref{eq:gevrey-rh} for every stated weight.
Its Mellin integral over $X\ge2$ is holomorphic for $\Re s>1/2$.
The bounded initial range adds an entire function. On $\Re s>1$,
absolute convergence and the substitution $u=\N\mathfrak n/X$
give
\[
 \int_0^\infty T_W(X)X^{-s}\,\frac{dX}{X}
 =-\frac{P'(s)}{P(s)}\widehat W(s).
\]
The integral is zero for all sufficiently small $X$, and
\eqref{eq:log-cutoff} gives its asserted continuation to
$\Re s>1/2$. If $\rho$ were a zero of $P$ there, choose
$J\subset(1,2)$ so short that the argument of $u^\rho$ varies
by less than $\pi/2$ on $J$. Lemma~\ref{lem:gevrey-bumps}
gives a nonzero nonnegative $W$ supported there. After a fixed
rotation the real part of $u^\rho$ is positive on the support,
so $\widehat W(\rho)\ne0$. The right side has a pole of
residue $-\operatorname{ord}_\rho(P)\widehat W(\rho)$, a
contradiction. Deleted Euler factors have zeros only on
$\Re s=0$, so the functional equation gives RH for the primitive
function.

For \eqref{eq:gevrey-separate}, the inner weight is
$W(u)(1-\chi(\lambda u))$ with $\lambda=L_X/\N\mathfrak c>1/2$.
It equals $W$ when $\lambda\ge2$; on $1/2<\lambda<2$ all
needed derivatives are uniformly bounded. The reciprocal-series
estimate therefore gives $O((X/\N\mathfrak c)^{\beta+\epsilon'})$
for any $0<\epsilon'<1-\beta$. Absolute summation over the outer
variable is bounded by
\[
 C X^{\beta+\epsilon'}L_X^{1-\beta-\epsilon'}\log(2L_X).
\]
For each requested $\epsilon>0$, choose
$0<\epsilon'\le\min(\epsilon,(1-\beta)/2)$ and enlarge the
constant. This yields \eqref{eq:gevrey-separate}.
\end{proof}

The improvement concerns the support of the remaining convolution.
It gives no estimate of the strength \eqref{eq:gevrey-rh}.
In particular, a logarithmic number of direct terms can still be
coupled to long inverse sums carrying every off-line zero pole.
The residue calculation below makes this last point independent of any
simplicity assumption on the zeros.

\section{A smaller cutoff from nearly analytic weights}
\label{sec:near-analytic}
We can shorten the direct range further while retaining a fixed,
nonnegative compactly supported test function. The needed construction
is an infinite convolution of interval densities. We give the details
so that no uniformity is inferred from an arbitrary smooth weight.

\begin{lemma}\label{lem:summable-widths}
Let $a_1,a_2,\ldots$ be positive real numbers with finite sum $S$.
There is a nonnegative $F\in C_c^\infty(\mathbb R)$ supported in
$[-S/2,S/2]$, of integral one, such that
\begin{equation}\label{eq:width-derivatives}
 \|F^{(k)}\|_\infty\le
 \frac{2^k}{a_1a_2\cdots a_{k+1}}\qquad(k\ge0).
\end{equation}
A translation places the support inside any interval of length greater
than $S$.
\end{lemma}
\begin{proof}
Let $b_j=a_j^{-1}\mathbf1_{[-a_j/2,a_j/2]}$. On the product of
these probability intervals, the sum of the coordinates converges
absolutely and uniformly because $\sum a_j<\infty$. Its distribution
is a probability measure $\mu$ supported in $[-S/2,S/2]$.
For every $n$, independence gives
$\mu=b_1*\cdots*b_n*\mu_{>n}$, where the last factor is a
probability measure.

The distributional derivative of $b_j$ is the difference of its two
endpoint masses divided by $a_j$, and has total variation $2/a_j$.
Fix $k\ge0$. Differentiate once in each of the first $k$ factors,
and leave the next two factors undifferentiated. The convolution
$b_{k+1}*b_{k+2}$ is continuous, bounded, and compactly supported,
with supremum at most $1/a_{k+1}$. Its convolution with the tail
measure is continuous with the same bound. Hence the $k$th
distributional derivative of $\mu$ is a continuous function bounded
by the right side of \eqref{eq:width-derivatives}. For $k=0$ this
identifies a continuous density $F$; the higher derivative identities
then give $F\in C^\infty$. Its nonnegativity, integral and support
follow from those of $\mu$.
\end{proof}

If the translated support lies in an interval inside $(0,\log2)$,
take $W(u)=F(\log u)$. The proof of
\eqref{eq:gevrey-factorial}, now using
\eqref{eq:width-derivatives} with $k=2A+4$, gives the general bound
\begin{equation}\label{eq:width-direct-bound}
 |A_W(Y)|\le C C_0^A Y^{-A}
       \prod_{j=1}^{2A+5}a_j^{-1}\qquad(A\ge1, Y\ge1).
\end{equation}
The constants $C,C_0$ are independent of $A$ and $Y$. All derivative
costs are printed in the product, rather than hidden in a constant
whose order can change with $Y$.

\begin{lemma}\label{lem:near-analytic-bumps}
For every $r>1$ and every nonempty open interval $J\subset(1,2)$,
there exists a nonzero nonnegative $W\in C_c^\infty(J)$ with
\begin{equation}\label{eq:near-analytic-derivatives}
 \|\mathcal D^kW\|_\infty
 \le C H^k k!\,[\log(e+k)]^{rk}\qquad(k\ge0).
\end{equation}
Every fixed weight satisfying this inequality obeys
\begin{equation}\label{eq:near-analytic-direct}
 |A_W(Y)|\le C'\exp\left(-\frac{c\sqrt Y}{[\log(3Y)]^r}\right)
 \qquad(Y\ge1).
\end{equation}
\end{lemma}
\begin{proof}
Choose
\[
 a_j=\frac{c_0}{j[\log(e+j)]^r}.
\]
The sum converges because $r>1$. Take $c_0$ sufficiently small and
translate the density of Lemma~\ref{lem:summable-widths} to have
support in $\log J$. Its derivative bound is at most
\[
 \frac{2^k}{c_0^{k+1}}(k+1)!
       [\log(e+k+1)]^{r(k+1)}.
\]
The ratio of this expression to $k![\log(e+k)]^{rk}$ is bounded
by a fixed constant to the power $k+1$. Indeed, the ratio of the
two logarithms raised to $rk$ is bounded, while the remaining
factor $(k+1)[\log(e+k+1)]^r$ has at most exponential growth.
This proves \eqref{eq:near-analytic-derivatives} for $W(u)=F(\log u)$.

For an arbitrary fixed weight satisfying
\eqref{eq:near-analytic-derivatives}, repeat the Mellin argument
with $k=2A+4$. It gives
\[
 |A_W(Y)|\le C_1 C_2^A Y^{-A}
       (2A+4)![\log(e+2A+4)]^{r(2A+4)}.
\]
Choose
\[
 A=\left\lfloor\frac{c_1\sqrt Y}{[\log(3Y)]^r}\right\rfloor.
\]
For all sufficiently large $Y$, $A\ge4$ and $k=2A+4\le3A$.
Also $\log(e+k)\le\log(3Y)$. Therefore
\[
 k[\log(e+k)]^r\le3c_1\sqrt Y.
\]
Using $k!\le k^k$, the last bound for $A_W$ is at most
$C_3Y^2(9c_1^2C_2)^A$, after enlarging a fixed constant.
Choose $c_1$ so small that $9c_1^2C_2<e^{-4}$.
Since $A/\log Y\to\infty$, the factor $Y^2$ is absorbed in
$e^A$ for large $Y$. This proves
\eqref{eq:near-analytic-direct}; bounded $Y$ is again absorbed.
\end{proof}

\subsection{Proof of the logarithmic cutoff criterion}
\begin{proof}[Proof of Theorem~\ref{thm:near-analytic-cutoff}]
For $\N\mathfrak d\le2D_X$, equation
\eqref{eq:near-analytic-direct} and the ideal count bound the
discarded part by
\[
 C X\exp\left(-\frac{c\sqrt{L_X/2}}
                         {[\log(3L_X/2)]^r}\right).
\]
The function $\sqrt Y/[\log(3Y)]^r$ is increasing for all sufficiently
large $Y$, which justifies using its lower endpoint. For large $X$,
$\log(3L_X/2)\le C_\kappa\ell_2(X)$, so the negative exponent has
magnitude at least
\[
 c'\ell_1(X)\ell_2(X)^{\kappa/2-r}.
\]
Since $\kappa/2-r>0$, the discarded part is $O_M(X^{-M})$
for every fixed $M$. The support and finite interchange arguments
are unchanged. The existence of localized nonnegative weights
proved above permits the same Mellin pole detection used in
Proposition~\ref{prop:gevrey-rh}, proving the RH equivalence.
The reciprocal-series estimate has the same uniformly bounded
inner-weight family and gives the displayed separate bound.
\end{proof}

For example, $r=3/2$ and $\kappa=4$ give direct norm at most
$2[\log(3X)]^2[\log(e+\log(3X))]^4$, with an error smaller
than every negative power of $X$. This range is smaller than
$(\log X)^{2+\epsilon}$ for every fixed $\epsilon>0$ once $X$
is large enough. The square-root estimate for its long inverse
sum remains unproved. No claim is made at $\kappa=2$ or for an
arbitrary compactly supported smooth test function.

\section{Why the zero poles remain}
We prove directly that the entire numerator does not cancel
the poles relevant to prime distribution.
For $H\ge1$ define the Dirichlet polynomial
\[
 A_H(s)=\sum_{\mathfrak c}
  a(\mathfrak c)\log\N\mathfrak c\,
  (\N\mathfrak c)^{-s}\chi(\N\mathfrak c/H).
\]

\begin{proposition}\label{prop:poles}
On every compact set of complex $s$, and for every $A>0$,
\[
 A_H(s)=-P'(s)+O_A(H^{-A}) \quad (H\longrightarrow\infty),
                                                               \tag{5.1}
\]
with the constant also depending on that compact set.
If $\rho$ is any zero of $P$ of multiplicity $r$, and $\Gamma$
is a sufficiently small positively oriented circle enclosing
$\rho$ and no other zero, then
\[
 \frac1{2\pi i}\int_\Gamma\frac{A_H(s)}{P(s)}\,ds
       \longrightarrow -r.                                  \tag{5.2}
\]
In particular, for all sufficiently large $H$,
$A_H/P$ has a nonremovable pole inside $\Gamma$.
\end{proposition}
\begin{proof}
The Mellin transform
$\widehat\chi(z)=\int_0^\infty\chi(u)u^{z-1}\,du$, initially
defined for $\Re z>0$, has a meromorphic continuation with
one simple pole at zero, of residue one. Indeed,
\[
 \widehat\chi(z)=-\frac1z\int_1^2\chi'(u)u^z\,du.
\]
The integral is entire and rapidly decreasing on each fixed
vertical strip by repeated integration by parts; the apparent
pole has the specified residue because $\int_1^2\chi'=-1$.
For a compact set of $s$, choose $c>0$ such that
$\Re(s+c)>1$ throughout that set. Mellin inversion gives
\[
 A_H(s)=\frac1{2\pi i}\int_{\Re z=c}
             -P'(s+z)\widehat\chi(z)H^z\,dz.
\]
Move the contour to $\Re z=-A$, crossing only the pole at zero.
Polynomial strip growth of $P'$ and the rapid decrease of
$\widehat\chi$ justify the move uniformly on the compact set.
The residue is $-P'(s)$ and the new integral is $O_A(H^{-A})$,
proving (5.1). The denominator is bounded away from zero on
$\Gamma$. Uniform convergence there therefore gives
\[
 \frac1{2\pi i}\int_\Gamma\frac{A_H(s)}{P(s)}\,ds
 \longrightarrow
 -\frac1{2\pi i}\int_\Gamma\frac{P'(s)}{P(s)}\,ds=-r.
\]
For large $H$ the integral is nonzero, so the quotient cannot
be holomorphic throughout the disk. This proves (5.2).
\end{proof}

For example, at a simple zero one has
$A_H(\rho)\to-P'(\rho)\ne0$, so the pole is visibly present
for all large $H$. The contour argument also handles multiple
zeros, without assuming that they are simple.

\section{Cancellation from rational content}\label{sec:content}

\subsection{The exact content identity}

We prove the local arithmetic input to
Theorem~\ref{thm:content-reduction}. Put
$\mathcal O=\mathbb Z[\omega]$ and $\eta=\chi_{-3}$,
and let $\mu$ be the ordinary integer M\"obius function.
The norm of a rounding error in the basis $1,\omega$, with
both coordinates of size at most $1/2$, is at most $3/4$.
Hence the Eisenstein norm is Euclidean, so every ideal is
principal. Its six units act by invertible integer matrices
on the coordinates; the rational content is therefore
independent of the generator.

For this cubic character,
\[
                        \psi((u))=1
        \qquad ((u,3\cdot269)=1,\ u\in\mathbb Z_{>0}).     \tag{6.1}
\]
In the splitting field $K(\sqrt[3]{269})$, complex conjugation
fixes the real cube root. It acts by inversion on the cyclic
Galois group over $K$.
The Frobenius conjugation relation gives
$\psi(\overline{\mathfrak a})=\psi(\mathfrak a)^{-1}$.
A rational ideal $(u)$ is fixed by conjugation.
Its character value consequently has both square and cube
equal to one, and is one.

Write $\alpha=u\alpha_0$ with
$\alpha_0=a+b\omega$ primitive, $\gcd(a,b)=1$, and
$q=\N\alpha_0=a^2-ab+b^2$. Away from the deleted primes,
\[
 \mu_K((\alpha))\psi((\alpha))
 =\mu(u)^2\eta(u)\mu(q)\psi((\alpha_0))
                       {\bf1}_{(u,q)=1}.                 \tag{6.2}
\]
To verify every zero condition, an inert rational prime dividing
$q$ would divide both coordinates of $\alpha_0$.
At a split rational prime, both conjugate prime ideals cannot
divide $(\alpha_0)$, since their product is the rational ideal.
The prime above three has been excluded. Thus the rational
prime valuations of $q$ are precisely the prime-ideal
valuations in $(\alpha_0)$, and
$\mu_K((\alpha_0))=\mu(q)$.
A squarefree rational content $u$ contributes two prime ideals
at each split prime and one at each inert prime, giving
$\mu_K((u))=\eta(u)$. A repeated rational prime in $u$,
or a common prime of $u$ and $q$, gives a repeated ideal factor
and a zero M\"obius value. This proves (6.2), using (6.1).

For $V\in C_c^\infty((1,2))$ define
\[
 M_{V,\chi}(Y,R)=\sum_{\mathfrak a}
   \mu_K(\mathfrak a)\psi(\mathfrak a)V(\N\mathfrak a/Y)
       (1-\chi(\operatorname{cont}(\mathfrak a)/R)).
\]
The exact lattice expression is
\begin{align*}
 M_{V,\chi}(Y,R)=\frac16
 \sum_{\substack{\alpha_0=a+b\omega\ne0\\(a,b)=1\\
                         (q,3\cdot269)=1}}
 &\mu(q)\psi((\alpha_0))
 \sum_{\substack{u\ge1\\(u,269q)=1}}\mu(u)^2\eta(u)\\
 &{}\times V(qu^2/Y)(1-\chi(u/R)).
\end{align*}
The factor $1/6$ accounts for associates; $\eta(u)$ already
vanishes at multiples of three.

\subsection{An elementary uniform estimate}

For $(h,3)=1$ and $T\ge1$, the squarefree expansion gives
\[
 \sum_{\substack{u\le T\\(u,h)=1}}\mu(u)^2\eta(u)
 =\sum_{\substack{d\le\sqrt T\\(d,3h)=1}}\mu(d)
   \sum_{e\mid\operatorname{rad}(h)}\mu(e)\eta(e)
       A\left(\frac{T}{d^2e}\right),\qquad
 A(v)=\sum_{1\le n\le v}\eta(n).
\]
Since $|A(v)|\le1$, the absolute value is at most
$\sqrt T\,\tau(h)$. Partial summation yields
\[
 \sum_{\substack{u\ge1\\(u,h)=1}}\mu(u)^2\eta(u)F(u/T)
       \ll_F\sqrt T\,\tau(h)                              \tag{6.3}
\]
for a bounded family of smooth functions $F$ with fixed compact
support in $(0,\infty)$.

Set
\[
 T_q=\sqrt{Y/q},\qquad
 F_{q,Y,R}(v)=V(v^2)(1-\chi((T_q/R)v)).
\]
Its support is fixed in $1<v<\sqrt2$.
If $T_q/R\ge2$ the cutoff factor is identically one there;
if $T_q/R\le1/\sqrt2$ it is identically zero.
All other parameter values lie in a fixed compact interval,
so every fixed derivative seminorm is uniformly bounded.
A nonzero term requires $q<2Y/R^2$.
The bounded interval $1/\sqrt2\le T_q<1$ can also be included
in (6.3), by increasing its constant.

The elementary lattice count is
$\#\{\alpha\in\mathcal O:0<\N\alpha\le Z\}\ll Z$.
Apply (6.3) with $h=269q$ and
$\tau(269q)\ll_\epsilon q^\epsilon$. For $Y\ge2$, $R\ge1$,
\[
 |M_{V,\chi}(Y,R)|
 \ll_\epsilon
 Y^{1/4}\sum_{\N\alpha_0<2Y/R^2}
                  (\N\alpha_0)^{-1/4+\epsilon}
 \ll_\epsilon Y^{1+\epsilon}R^{-3/2}.                    \tag{6.4}
\]
If $2Y/R^2<1$, the sum is empty.
This proof is unconditional.

\subsection{A smooth refinement under zeta RH}

Assume classical RH for $\zeta$ only in this subsection.
The scalar series in (6.3), initially for $\Re s>1$, is
\[
 F_h(s)=\sum_{\substack{u\ge1\\(u,h)=1}}
       \frac{\mu(u)^2\eta(u)}{u^s}
 =\frac{L(s,\eta)}
        {\zeta(2s)(1-3^{-2s})}
       \prod_{p\mid h}(1+\eta(p)p^{-s})^{-1}.              \tag{6.5}
\]
The factor at three is included because its squarefree Euler
factor must be one. The final product accounts for all excluded
prime divisors of $h=269q$.

For fixed $\sigma_0,\epsilon>0$,
\[
 \sup_{\Re s\ge\sigma_0}
   \prod_{p\mid h}|1+\eta(p)p^{-s}|^{-1}
       \ll_{\sigma_0,\epsilon}h^\epsilon.                 \tag{6.6}
\]
Indeed each factor is at most $(1-p^{-\sigma_0})^{-1}$.
For sufficiently large $p$ its logarithm is at most
$\epsilon\log p$, and the remaining finitely many factors
give a constant. This proves uniformity even when $q$ grows.

Under zeta RH, (6.5) is holomorphic on $\Re s>1/4$.
The zeta pole at $2s=1$ produces a zero, not a pole.
The finite Euler denominators have zeros only on $\Re s=0$.
The reciprocal $1/\zeta(z)$ has polynomial growth on every
closed inner strip of $\Re z>1/2$: remove the pole by using
$(z-1)\zeta(z)$, choose an analytic logarithm in that zero-free
half-plane, anchor it at $2+it$ by the Euler product, and
apply Borel--Carath\'eodory on disks contained in a slightly
larger inner strip. The functional equation and standard
strip bounds give the upper real part
$O(\log(3+|t|))$ needed in that argument.
The numerator $L(s,\eta)$ has polynomial strip growth as well.

Mellin inversion, rapid Mellin decay, and displacement to
$\Re s=1/4+\delta$ therefore imply
\[
 \sum_{\substack{u\ge1\\(u,h)=1}}\mu(u)^2\eta(u)F(u/T)
       \ll_{\delta,\epsilon,F}T^{1/4+\delta}h^\epsilon,
             \qquad \delta,\epsilon>0.                  \tag{6.7}
\]
There are no crossed poles. All constants involve only
finitely many smooth seminorms, so (6.7) applies uniformly
to $F_{q,Y,R}$. In particular, no sharp content threshold
has replaced the smooth cutoff.

Choose $0<\delta<1/4$. Writing $\theta=1/4+\delta$ and
counting primitive vectors gives
\[
 |M_{V,\chi}(Y,R)|
 \ll_{\delta,\epsilon}
    Y^{\theta/2}(Y/R^2)^{1-\theta/2+\epsilon}
 =Y^{1+\epsilon}R^{-7/4+\delta-2\epsilon}.
\]
A nonzero sum has $R\le\sqrt{2Y}$. Choose internal positive
$\delta,\epsilon$ smaller than the desired final error
exponent, and absorb $R^\delta$ into $Y^{\delta/2}$.
Thus for every final $\epsilon>0$,
\[
                M_{V,\chi}(Y,R)
                    \ll_\epsilon Y^{1+\epsilon}R^{-7/4}.
                                                               \tag{6.8}
\]

\subsection{Proof of the content reduction theorem}

\begin{proof}[Proof of Theorem~\ref{thm:content-reduction}]
For each direct ideal $\mathfrak c$, the inner weight is
\[
 V_{\mathfrak c}(v)=W(v)(1-\chi((L_X/\N\mathfrak c)v)).
\]
It ranges over a bounded smooth family. Indeed its varying
parameter is in $[1/2,2]$ unless the subtracted cutoff
is identically zero.
Apply (6.4) at $R=Y_{\mathfrak c}^{1/3}$, or (6.8) at
$R=Y_{\mathfrak c}^{2/7}$ under zeta RH.
The discarded inner content range is
$O_\epsilon(Y_{\mathfrak c}^{1/2+\epsilon})$.
Summation over the direct ideals costs
\[
 \sum_{\N\mathfrak c<2L_X}\log\N\mathfrak c\,
                         Y_{\mathfrak c}^{1/2+\epsilon}
 \ll_\epsilon X^{1/2+\epsilon}L_X^{1/2}\log(2L_X).
\]
The last factors are absorbed by an arbitrarily small increase
in the exponent of $X$.
Theorem~\ref{thm:near-analytic-cutoff} supplies the other
discarded range, with arbitrarily rapid decay.
This proves the two approximation identities.
The RH equivalences follow from the equivalence in that
theorem, because the new error is at the square-root scale
for every positive error exponent.

In the retained sum, the inner supports give
$Y_{\mathfrak c}<\N\mathfrak d<2Y_{\mathfrak c}$
and $\operatorname{cont}(\mathfrak d)<2Y_{\mathfrak c}^\gamma$.
Dividing the norm by the square of the content gives
$\N\alpha_0>\tfrac14Y_{\mathfrak c}^{1-2\gamma}$.
All estimates hold for sufficiently large $X$; the initial
bounded range is absorbed into the constants.
\end{proof}

\subsection{The primitive inverse range remains}

The ordinary M\"obius factor in (6.2) is
$\mu(a^2-ab+b^2)$, and it has not been removed.
More precisely, let
\[
 h(n)=\sum_{\substack{\N\mathfrak a=n\\
                     \operatorname{cont}(\mathfrak a)=1}}
                         \mu_K(\mathfrak a)\psi(\mathfrak a).
\]
It is multiplicative, supported on squarefree integers with
all prime factors $1\pmod3$, and
\[
 h(p)=-t(p),\qquad
 t(p)=\#\{z\bmod p:z^3=269\}-1\in\{2,-1\}.
\]
At an inert prime there is no primitive ideal; at a split
prime the two conjugate ideals give the stated trace.
Initially on $\Re s>1$,
\[
 \sum_n h(n)n^{-s}=\frac{C(s)}{P(s)},\qquad
 C(s)=
 \prod_{\substack{p\equiv2\ (3)\\p\ne269}}(1-p^{-2s})^{-1}
 \prod_{p\equiv1\ (3)}
       \frac{1-t(p)p^{-s}}{1-t(p)p^{-s}+p^{-2s}}.           \tag{6.9}
\]
The product for $C$ is absolutely convergent and nowhere zero
on $\Re s>1/2$. Its factors differ from one by
$O(p^{-2\Re s})$. The split denominators are products
$(1-zp^{-s})(1-\bar z p^{-s})$, $|z|=1$.
The numerator $1+p^{-s}$ has no zero there; for a numerator
$1-2p^{-s}$ use $p\ge7$ to get $2p^{-\Re s}<1$.
Thus every possible zero pole of $P$ in $\Re s>1/2$
persists in the primitive coefficient series.

This is also a literal part of the restricted convolution.
For $\mathfrak c=(2)$, of norm four and character one,
and sufficiently large $X$, its primitive inner portion is
\[
                         (\log4)\sum_n h(n)W(4n/X).
\]
Other direct terms could cancel it. No lower bound for the
whole remainder is inferred, but the long primitive component
has not been estimated at the square-root scale.
The content deletion is a strict support reduction, not an
improved exponent for the complete prime sum.


\section*{What remains open}
For the explicitly constructed nearly analytic weights, all inverse variables up to $X/L_X$ have been discarded with arbitrarily rapid decay, where $L_X=(\log(3X))^2[\log(e+\log(3X))]^\kappa$ and any fixed $\kappa>2$ is permitted. The remaining direct norm is below $2L_X$. Rational contents beyond the smooth threshold $Y^{1/3}$ in each inverse window can also be removed at square-root cost; under zeta RH the threshold is $Y^{2/7}$. The retained primitive norm still ranges up to $2Y$, and the coefficients contain $\mu(a^2-ab+b^2)\psi((a+b\omega))$. A square-root estimate for the retained convolution, for every localized nonnegative weight in the stated class, is equivalent to the cubic Riemann hypothesis.

Separate estimates from the named zero-free theorem yield exponent $B_*+\epsilon$, with logarithmic factors; the complete explicit formula gives $B_*$. The smaller support does not remove zero poles, as the contour calculation proves. No cancellation among the remaining outer terms at the square-root scale has been established. The cubic Riemann hypothesis remains open.

\section*{Method and availability}
The mathematics and text were produced by GPT-6 Astra in ChatGPT Pro mode.
The cited Hecke functional equation, local zero count, and classical
derivative classes are the background inputs. The cutoff, content,
Mellin-continuation, and retained-pole arguments are proved here.
The named recent zero-free estimates are used only for comparison of
exponents and are not re-proved. The accompanying records distinguish
statement-only reconstructions from targeted checks by the same model.
No human review or formal verification is claimed.
The original text is dedicated to the public domain under CC0 1.0;
cited works retain their own terms. Sources and verification records are
at \url{\workbenchurl}, with version records at \url{\recordurl}.

\begin{thebibliography}{99}
\bibitem{Poonen} Poonen, B. (2015).
\emph{Tate's thesis}. Theorems 5.16 and 5.22(b,c), pp.~34--39.
\url{https://math.mit.edu/~poonen/786/notes.pdf}.
\bibitem{OpenAIreflection} OpenAI. (2026a, September 25).
\emph{An unconditional first moment for cubic Gauss sums}.
Lemma 4.2, p.~14.
\href{https://github.com/openai/math/blob/main/preprints/An-unconditional-first-moment-for-cubic-Gauss-sums-September-25-2026/paper.pdf}{Dated primary release}.
\bibitem{OpenAIzero} OpenAI. (2026b, September 30).
\emph{The quasi-Riemann hypothesis: A zero-free half-plane
$\Re s>7/8$}. Theorem 1.1, p.~4.
\href{https://github.com/openai/math/blob/main/preprints/The-Quasi-Riemann-Hypothesis-September-30-2026/paper.pdf}{Dated primary release}.
\bibitem{Bramble} bramble\_trace. (2026, October 26).
\emph{Mixed product moments and natural zero density for cubic
Hecke characters} (Version 5). Section 2.2.
\href{https://rstation.ai/problems/riemann_hypothesis/papers/52}{Dated rStation paper, version 5}.
\bibitem{Gevrey} Gevrey, M. (1918). Sur la nature analytique des solutions des \'equations aux d\'eriv\'ees partielles. Premier m\'emoire. \emph{Annales scientifiques de l'\'Ecole Normale Sup\'erieure, 35}, 129--190. Part I, Section 1, p.~132. \url{https://doi.org/10.24033/asens.706}.
\bibitem{Liu} Liu, B. (2026, November 13).
\emph{Joint marked moments and optimized zero-free half-planes
for unitary Hecke $L$-functions} (Version 2).
Theorems 1.2 and 8.29, pp.~3 and 78--80; equations (42) and (44), pp.~27--29.
\url{https://hexagonmath.org/2611.00421v2}.
\bibitem{GM} Greni\'e, L., \& Molteni, G. (2019).
\emph{An effective Chebotarev density theorem under GRH}
(Version 3, May 28). Appendix A, Theorem A.1, is unconditional.
\url{https://arxiv.org/abs/1709.07609v3}.
\end{thebibliography}
\end{document}
